% part: intuitionistic-logic
% chapter: propositions-as-types
% section: terms-to-proofs

\documentclass[../../../include/open-logic-section]{subfiles}

\begin{document}

\olfileid{int}{pty}{tp}

\olsection{প্রমাণপদ থেকে \usetoken{P}{derivation} পুনর্গঠন}

এবার উলটো দিকটি দেখি: একটি প্রমাণপদ~$N$ থেকে
\Log{N2i}-তে !!a{proof} পুনর্গঠন করি। সাধারণভাবে,
দেওয়া প্রমাণপদে মুক্তভাবে ঘটা প্রত্যেক
চলরাশি~$x$-এর জন্য $x : !A$ যুগল আছে এমন
প্রসঙ্গের সাপেক্ষেই এটি সম্ভব। তবে প্রথমে
মুক্ত চলরাশিহীন একটি উদাহরণ নিই, এই পদটি:
\[
\lambd[\typeof{x}{!A \land
  !B}][\pair{\proj{2}{x}}{\proj{1}{x}}]
\]
এটি $\lambd[\typeof{x}{!A \land !B}][N_1]$
রূপের। \Log{tN2i}-এর বিধিগুলির মধ্যে কেবল
\Intro{\lif}-ই এই রূপের পদ দেয়। তাই $N$-এ
সংকেতায়িত !!{proof}-এর শেষ বিধি
\Intro{\lif}; অর্থাৎ এর শেষাংশ
  \[
  \Axiom$x : !A\land !B \fCenter \pair{\proj{2}{x}}{\proj{1}{x}} :
    !C$
  \RightLabel{\Intro\lif}
  \UnaryInf$\fCenter \lambd[\typeof{x}{!A \land
  !B}][\pair{\proj{2}{x}}{\proj{1}{x}}] : (!A \land !B) \lif !C$
  \DisplayProof
  \]
$!C$ কী, তা এখনও জানি না; কিন্তু পূর্বপক্ষ
$!A \land !B$ এবং পূর্বধারণার প্রসঙ্গে
$x:!A \land !B$ থাকতে হবে, তা জানি।

এবার পূর্বধারণার প্রমাণপদটি নির্ধারিত:
$\pair{\proj{2}{x}}{\proj{1}{x}}$।
\emph{এই পদটি} যে !!{proof} সংকেতায়িত করে,
তার শেষ বিধি \Intro{\land}। $!C$ এখনও
অজানা, কিন্তু \Intro{\land}-এর সিদ্ধান্ত
হওয়ায় এটি একটি যৌগ হতে হবে, অর্থাৎ
$!C \ident (!C_1 \land !C_2)$।
পুনর্গঠন চালিয়ে পাই:
  \[
  \Axiom$x : !A\land !B \fCenter \proj{2}{x} : !C_1$
  \Axiom$x : !A\land !B \fCenter \proj{1}{x} : !C_2$
  \RightLabel{\Intro\land}
  \BinaryInf$x : !A\land !B \fCenter \pair{\proj{2}{x}}{\proj{1}{x}} :
    !C_1 \land !C_2$
  \RightLabel{\Intro\lif}
  \UnaryInf$\fCenter \lambd[\typeof{x}{!A \land
  !B}][\pair{\proj{2}{x}}{\proj{1}{x}}] : (!A \land !B) \lif (!C_1 \land !C_2)$
  \DisplayProof
  \]
প্রমাণপদ $\pi_2(x)$ দেখায়, বাঁ দিকের উপরের
সিকোয়েন্টটি \Elim{\land}[2] দিয়ে পাওয়া;
অর্থাৎ তার পূর্বধারণা $!D \land !C_1$
রূপের। ডান দিকে $!C_2$-তে পৌঁছানো অনুমানটি
\Elim{\land}[1] হতে হবে, আর তার পূর্বধারণা
$!C_2 \land !E$ রূপের।
  \[
  \Axiom$x : !A\land !B \fCenter x: !D\land !C_1$
  \RightLabel{\Elim\land[2]}
  \UnaryInf$x : !A\land !B \fCenter \proj{2}{x} : !C_1$
  \Axiom$x : !A\land !B \fCenter x: !C_2\land !E$
  \RightLabel{\Elim\land[1]}
  \UnaryInf$x : !A\land !B \fCenter \proj{1}{x} : !C_2$
  \RightLabel{\Intro\land}
  \BinaryInf$x : !A\land !B \fCenter \pair{\proj{2}{x}}{\proj{1}{x}} :
    !C_1 \land !C_2$
  \RightLabel{\Intro\lif}
  \UnaryInf$\fCenter \lambd[\typeof{x}{!A \land
  !B}][\pair{\proj{2}{x}}{\proj{1}{x}}] : (!A \land !B) \lif (!C_1 \land !C_2)$
  \DisplayProof
  \]
উপরের দুই সিকোয়েন্টের প্রমাণপদ এখন
চলরাশি; তাই সেগুলি স্বতঃসিদ্ধ হতে হবে।
এতে $!C_1$, $!C_2$, $!D$ ও $!E$ নির্ধারিত
হয়: $!D \ident !A$ এবং $!C_1 \ident !B$,
নইলে বাঁ সিকোয়েন্ট স্বতঃসিদ্ধ নয়; আর
$!C_2 \ident !A$ এবং $!E \ident !B$,
নইলে ডান সিকোয়েন্ট স্বতঃসিদ্ধ নয়।
এবার সম্পূর্ণ প্রমাণটি ফিরে পেলাম:
  \[
  \Axiom$x : !A\land !B \fCenter x: !A\land !B$
  \RightLabel{\Elim\land[2]}
  \UnaryInf$x : !A\land !B \fCenter \proj{2}{x} : !B$
  \Axiom$x : !A\land !B \fCenter x: !A\land !B$
  \RightLabel{\Elim\land[1]}
  \UnaryInf$x : !A\land !B \fCenter \proj{1}{x} : !A$
  \RightLabel{\Intro\land}
  \BinaryInf$x : !A\land !B \fCenter \pair{\proj{2}{x}}{\proj{1}{x}} :
    !B \land !A$
  \RightLabel{\Intro\lif}
  \UnaryInf$\fCenter \lambd[\typeof{x}{!A \land
  !B}][\pair{\proj{2}{x}}{\proj{1}{x}}] : (!A \land !B) \lif (!B \land !A)$
  \DisplayProof
  \]


\end{document}
